%%%%%%%%%%%%%%%%%%%%%%%%%%%%%%%%%%%%%%%%%%%%%%%%%%%%%%%%%%%%%%%%%%%%%%
% Lớp 11 | Hàm số mũ và hàm số lôgarít | [Giảm] Các phương pháp giải được giảm tải
% Dạng 1D6?5-2: Phương pháp lôgarít hóa, mũ cho PT mũ, lôgarít
%%%%%%%%%%%%%%%%%%%%%%%%%%%%%%%%%%%%%%%%%%%%%%%%%%%%%%%%%%%%%%%%%%%%%%


% nguồn: 2026-11-15_11_1D6_Mu_Logarit (nhập 14/11/2026)
\begin{bt}%[1D6H5-2]%[ma:90274345]
% nguon: Archive/SP-4Phan-Dot3/Data/CD_1_OTHK2_De03.tex #4411

	Tìm tập hợp nghiệm của bất phương trình $3^x-2\cdot5^x<0$.
	\loigiai{
		Ta có
	\begin{eqnarray*}
		 3^x-2\cdot5^x<0
		 &\Leftrightarrow&
		 3^x<2\cdot5^x\\
		  &\Leftrightarrow&
		 \dfrac{3^x}{5^x}<2\\
		  &\Leftrightarrow&
		 \left(\dfrac{3}{5}\right)^x<2\\
		 &\Leftrightarrow&
		 x>\log_{\tfrac{3}{5}} 2.
	\end{eqnarray*}
	Vậy tập hợp nghiệm của bất phương trình $3^x-2\cdot5^x<0$ là $S=\left(\log_{\tfrac{3}{5}} 2; +\infty \right)$.
	}
\end{bt}
